\documentclass[10pt,a4paper]{amsart}
\usepackage{amsmath,amssymb,amsthm,graphicx}
% Retained preamble for the Stacks Project chapter {CH} (shared project preamble).
% The chapter file pinned-original-source.tex only carries \input{{preamble}} on its first line,
% so no preamble text lives inside the chapter itself.  The Stacks Project uses one single shared
% preamble for every chapter: the snapshot stages it three times (sources/stacks-chow/preamble.tex,
% sources/stacks-etale-cohomology/preamble.tex and focused-seeds/AGP/sources/stacks-snapshot-20261011/
% preamble.tex) and those copies are identical in size.  The block below is copied mechanically,
% line for line, from sources/stacks-chow/preamble.tex lines 14-18 and 154-217: the Xy-pic setup and
% the native theorem declarations together with the mathematical macros.  Nothing was added, removed
% or redefined; the wrapper class and the amsmath/amssymb/amsthm/graphicx base packages are supplied
% by the assembler.
% For commutative diagrams we use Xy-pic
% Chapter: stacks-spaces-more-morphisms.tex -- More on Morphisms of Spaces.
\usepackage[all]{xy}

% We use 2cell for 2-commutative diagrams.
\xyoption{2cell}

% Theorem environments.
%
\theoremstyle{plain}
\newtheorem{theorem}[subsection]{Theorem}
\newtheorem{proposition}[subsection]{Proposition}
\newtheorem{lemma}[subsection]{Lemma}

\theoremstyle{definition}
\newtheorem{definition}[subsection]{Definition}
\newtheorem{example}[subsection]{Example}
\newtheorem{exercise}[subsection]{Exercise}
\newtheorem{situation}[subsection]{Situation}

\theoremstyle{remark}
\newtheorem{remark}[subsection]{Remark}
\newtheorem{remarks}[subsection]{Remarks}

\numberwithin{equation}{subsection}

% Macros
%
\def\lim{\mathop{\mathrm{lim}}\nolimits}
\def\colim{\mathop{\mathrm{colim}}\nolimits}
\def\Spec{\mathop{\mathrm{Spec}}}
\def\Hom{\mathop{\mathrm{Hom}}\nolimits}
\def\Ext{\mathop{\mathrm{Ext}}\nolimits}
\def\SheafHom{\mathop{\mathcal{H}\!\mathit{om}}\nolimits}
\def\SheafExt{\mathop{\mathcal{E}\!\mathit{xt}}\nolimits}
\def\Sch{\mathit{Sch}}
\def\Mor{\mathop{\mathrm{Mor}}\nolimits}
\def\Ob{\mathop{\mathrm{Ob}}\nolimits}
\def\Sh{\mathop{\mathit{Sh}}\nolimits}
\def\NL{\mathop{N\!L}\nolimits}
\def\CH{\mathop{\mathrm{CH}}\nolimits}
\def\proetale{{pro\text{-}\acute{e}tale}}
\def\etale{{\acute{e}tale}}
\def\QCoh{\mathit{QCoh}}
\def\Ker{\mathop{\mathrm{Ker}}}
\def\Im{\mathop{\mathrm{Im}}}
\def\Coker{\mathop{\mathrm{Coker}}}
\def\Coim{\mathop{\mathrm{Coim}}}

% Boxtimes
%
\DeclareMathSymbol{\boxtimes}{\mathbin}{AMSa}{"02}

%
% Macros for moduli stacks/spaces
%
\def\QCohstack{\mathcal{QC}\!\mathit{oh}}
\def\Cohstack{\mathcal{C}\!\mathit{oh}}
\def\Spacesstack{\mathcal{S}\!\mathit{paces}}
\def\Quotfunctor{\mathrm{Quot}}
\def\Hilbfunctor{\mathrm{Hilb}}
\def\Curvesstack{\mathcal{C}\!\mathit{urves}}
\def\Polarizedstack{\mathcal{P}\!\mathit{olarized}}
\def\Complexesstack{\mathcal{C}\!\mathit{omplexes}}
% \Pic is the operator that assigns to X its picard group, usage \Pic(X)
% \Picardstack_{X/B} denotes the Picard stack of X over B
% \Picardfunctor_{X/B} denotes the Picard functor of X over B
\def\Pic{\mathop{\mathrm{Pic}}\nolimits}
\def\Picardstack{\mathcal{P}\!\mathit{ic}}
\def\Picardfunctor{\mathrm{Pic}}
\def\Deformationcategory{\mathcal{D}\!\mathit{ef}}
\usepackage{cleveref}
\crefname{originalreference}{original reference}{original references}
\Crefname{originalreference}{Original reference}{Original references}
\makeatletter
\newlabel{decent-spaces-remark-recall}{{03II}{source}{}{sourcepointer-1}{}}
\newlabel{decent-spaces-remark-recall@cref}{{[originalreference][0][]03II}{[1][0][]source}{}{}{}}
\makeatother
\crefname{originalreference}{original reference}{original references}
\Crefname{originalreference}{Original reference}{Original references}
\makeatletter
\newlabel{more-algebra-lemma-fibre-product-finite-type}{{00IT}{source}{}{sourcepointer-2}{}}
\newlabel{more-algebra-lemma-fibre-product-finite-type@cref}{{[originalreference][0][]00IT}{[1][0][]source}{}{}{}}
\makeatother
\crefname{originalreference}{original reference}{original references}
\Crefname{originalreference}{Original reference}{Original references}
\makeatletter
\newlabel{more-algebra-lemma-points-of-fibre-product}{{0B7J}{source}{}{sourcepointer-3}{}}
\newlabel{more-algebra-lemma-points-of-fibre-product@cref}{{[originalreference][0][]0B7J}{[1][0][]source}{}{}{}}
\makeatother
\crefname{originalreference}{original reference}{original references}
\Crefname{originalreference}{Original reference}{Original references}
\makeatletter
\newlabel{more-algebra-situation-module-over-fibre-product}{{08KH}{source}{}{sourcepointer-4}{}}
\newlabel{more-algebra-situation-module-over-fibre-product@cref}{{[originalreference][0][]08KH}{[1][0][]source}{}{}{}}
\makeatother
\crefname{originalreference}{original reference}{original references}
\Crefname{originalreference}{Original reference}{Original references}
\makeatletter
\newlabel{morphisms-lemma-affine-equivalence-algebras}{{01SA}{source}{}{sourcepointer-5}{}}
\newlabel{morphisms-lemma-affine-equivalence-algebras@cref}{{[originalreference][0][]01SA}{[1][0][]source}{}{}{}}
\makeatother
\crefname{originalreference}{original reference}{original references}
\Crefname{originalreference}{Original reference}{Original references}
\makeatletter
\newlabel{spaces-morphisms-lemma-pushforward}{{03M9}{source}{}{sourcepointer-6}{}}
\newlabel{spaces-morphisms-lemma-pushforward@cref}{{[originalreference][0][]03M9}{[1][0][]source}{}{}{}}
\makeatother
\title{Glueing an integral morphism so that over a point the fibre becomes a single point with prescribed residue field}
\author{The Stacks Project Authors}
\date{stacks-project:spaces-more-morphisms.tex:local-snapshot-20261011; GNU FDL 1.2 or later}
\begin{document}
\maketitle
\noindent\textit{Editorial note.} Counted span is the author's own proof environment, source lines 7289-7346, immediately following the statement at 7273-7287.\par\medskip
\noindent\textit{Editorial note.} Six cross-chapter pointers (Decent Spaces, More on Algebra, Morphisms, Morphisms of Spaces) are recorded in assembly.external\_references.\par\medskip
\noindent\textit{Editorial note.} The final finiteness statement is proved inside the span from the fibre product algebra being finitely generated over R.\par\medskip
\noindent\textit{Allowed standard prerequisites (editorial).}\begin{itemize}
\item Fibre products of algebras and their points. More on Algebra, situation more-algebra-situation-module-over-fibre-product and lemmas more-algebra-lemma-points-of-fibre-product and more-algebra-lemma-fibre-product-finite-type, cited throughout. Describes the fibre product algebra C' and its points, which is how the new fibre X'\_y is controlled.
\item Affine equivalence for algebras and modules. Morphisms, lemma morphisms-lemma-affine-equivalence-algebras with Morphisms of Spaces, lemma spaces-morphisms-lemma-pushforward, cited in the proof. Moves the construction between quasi-coherent sheaves and their global sections.
\item Recall on decent spaces. Decent Spaces, remark decent-spaces-remark-recall, cited in the proof. Allows the residue field of the chosen point to be treated as a finite separable extension in the decent case.
\item Integral morphisms are affine and closed. standard Stacks facts about integral morphisms used throughout the proof. Keeps the constructed factorization X -> X' -> Y integral over Y and lets the fibre condition be checked affine-locally.
\end{itemize}\medskip
\ifdefined\hypertarget\hypertarget{sourcepointer-1}{}\fi
\noindent\textit{Source reference.} Original reference 03II: focused-seeds/geometry/S9-spaces-construct-glueing/ref-tags.txt, source lines 4--4. Navigation pointer only: the counted proof cites this statement by \textbackslash{}ref as a standing input; the deduction the counted proof has to make is carried out inside the counted span, and no part of that deduction is delegated to this pointer.\par\medskip
\ifdefined\hypertarget\hypertarget{sourcepointer-2}{}\fi
\noindent\textit{Source reference.} Original reference 00IT: focused-seeds/geometry/S9-spaces-construct-glueing/ref-tags.txt, source lines 5--5. Navigation pointer only: the counted proof cites this statement by \textbackslash{}ref as a standing input; the deduction the counted proof has to make is carried out inside the counted span, and no part of that deduction is delegated to this pointer.\par\medskip
\ifdefined\hypertarget\hypertarget{sourcepointer-3}{}\fi
\noindent\textit{Source reference.} Original reference 0B7J: focused-seeds/geometry/S9-spaces-construct-glueing/ref-tags.txt, source lines 6--6. Navigation pointer only: the counted proof cites this statement by \textbackslash{}ref as a standing input; the deduction the counted proof has to make is carried out inside the counted span, and no part of that deduction is delegated to this pointer.\par\medskip
\ifdefined\hypertarget\hypertarget{sourcepointer-4}{}\fi
\noindent\textit{Source reference.} Original reference 08KH: focused-seeds/geometry/S9-spaces-construct-glueing/ref-tags.txt, source lines 7--7. Navigation pointer only: the counted proof cites this statement by \textbackslash{}ref as a standing input; the deduction the counted proof has to make is carried out inside the counted span, and no part of that deduction is delegated to this pointer.\par\medskip
\ifdefined\hypertarget\hypertarget{sourcepointer-5}{}\fi
\noindent\textit{Source reference.} Original reference 01SA: focused-seeds/geometry/S9-spaces-construct-glueing/ref-tags.txt, source lines 8--8. Navigation pointer only: the counted proof cites this statement by \textbackslash{}ref as a standing input; the deduction the counted proof has to make is carried out inside the counted span, and no part of that deduction is delegated to this pointer.\par\medskip
\ifdefined\hypertarget\hypertarget{sourcepointer-6}{}\fi
\noindent\textit{Source reference.} Original reference 03M9: focused-seeds/geometry/S9-spaces-construct-glueing/ref-tags.txt, source lines 9--9. Navigation pointer only: the counted proof cites this statement by \textbackslash{}ref as a standing input; the deduction the counted proof has to make is carried out inside the counted span, and no part of that deduction is delegated to this pointer.\par\medskip
\begin{lemma}
\label{lemma-construct-glueing}
Let $S$ be a scheme. Let $f : X \to Y$ be an integral morphism
of algebraic spaces over $S$.
Let $y \in |Y|$ be a point which can be represented by a closed immersion
$y : \Spec(k) \to Y$. Then there exists
a factorization $X \to X' \to Y$ of $f$ such that
\begin{enumerate}
\item $X' \to Y$ is integral,
\item $X \to X'$ is an isomorphism over $X' \setminus X'_y$,
\item $X'_y$ has a unique point $x'$ with $\kappa(x') = k$.
\end{enumerate}
Moreover, if $f$ is finite and $Y$ is locally Noetherian, then
$X' \to Y$ is finite.
\end{lemma}

\begin{proof}
By Morphisms of Spaces, Lemma \ref{spaces-morphisms-lemma-pushforward}
the sheaves $f_*\mathcal{O}_X$, $(X_y \to Y)_*\mathcal{O}_{X_y}$, and
$y_*\mathcal{O}_{\Spec(k)}$ are quasi-coherent sheaves of
$\mathcal{O}_Y$-algebras. Consider the maps
$$
f_*\mathcal{O}_Y \longrightarrow
(X_y \to Y)_*\mathcal{O}_{X_y} \longleftarrow
y_*\mathcal{O}_{\Spec(k)}
$$
The fibre product is a quasi-coherent sheaf of $\mathcal{O}_Y$-algebras
$\mathcal{A}'$ and we can define $X' \to Y$ as the relative spectrum
of $\mathcal{A}'$ over $Y$, see
Morphisms, Lemma \ref{morphisms-lemma-affine-equivalence-algebras}.
This construction commutes with arbitrary change of base.
In particular, it is clear that over the open subspace
$|Y| \setminus \{y\}$ the morphism $X \to X'$ is an isomorphism
and over $|Y| \setminus \{y\}$ the morphism $X' \to Y$ is integral.
It remains to prove the statements in a small neighbourhood of $y$.
Choose an affine scheme $V = \Spec(R)$ and an \'etale
morphism $\varphi : V \to Y$ such that $y$ is in the image of
$\varphi$. Then $V_y$ is a closed subscheme of $V$
\'etale over $k$, whence consists of finitely many points
each with residue field separable over $k$
(see Decent Spaces, Remark \ref{decent-spaces-remark-recall}).
After shrinking $V$
we may assume there is a unique closed point $v = \Spec(l) \to V$
mapping to $y$ with $l/k$ finite separable.
We may write $V \times_Y X = \Spec(C)$ with $R \to C$
an integral ring map. The stated compatibility with
base change gives us that $U \times_X Y' = \Spec(C')$ where
$$
C' = C \times_{C \otimes_R l} l
$$
Since $R \to l$ is surjective, also $C \to C \otimes_R l$
is surjective and we see that this is a fibre product of the
kind studied in More on Algebra, Situation
\ref{more-algebra-situation-module-over-fibre-product}
(with $A, A', B, B'$ corresponding to $C \otimes_R l, C, l, C'$).
Observe that $C'$ is an $R$-subalgebra of $C$ and hence
is integral over $R$; this proves (1).
Finally, More on Algebra, Lemma
\ref{more-algebra-lemma-points-of-fibre-product}
shows that $V \times_X Y' = \Spec(C')$
has a unique point $y''$ lying over $v$ with residue $l$
(this corresponds with the obvious surjective map $C' \to l$).
Thus $X_y \times_{\Spec(k)} \Spec(l)$ has a unique point
with residue field $l$. Since $l/k$ is finite separable,
this implies $X'_y$ has a unique point with residue field $k$, i.e.,
(3) holds.

\medskip\noindent
To prove the final statement, observe that if $Y$ is locally Noetherian,
then $R$ is a Noetherian ring and if $f$ is finite, then $R \to C$ is
finite. Then $C'$ is a finite type $R$-algebra by More on Algebra,
Lemma \ref{more-algebra-lemma-fibre-product-finite-type}.
This proves that $X' \to Y$ is finite.
\end{proof}

\end{document}
